Mula-mula terdapat $1$ kotak yang berisi $50^{25}$ batu. Perhatikan operasi berikut yang dilakukan secara berulang:

Ambil $1$ batu dari setiap kotak yang berisi paling sedikit $1$ batu. Lalu, siapkan sebuah kotak kosong baru dan masukkan semua batu yang tadi diambil ke dalam kotak baru tersebut.

Setelah operasi ke-$50^{25}$ selesai, ada berapa banyak kotak yang berisi paling sedikit $1$ batu?
\begin{solusi}
\section{Motivasi dan Walkthrough}

\begin{itemize}
    \item Operasi ini sebenarnya punya nama: \textcrimson{\textbf{Bulgarian Solitaire}}. Teknik ini pertama kali dibahas di 1983 Scientific American column oleh Martin Gardner.
    \item Coba main-main dulu sama angka kecil, misal mulai dari $6$ batu. Tulis banyaknya batu tiap kotak, urut dari yang terbanyak, terus jalanin operasinya beberapa langkah. Nyadar sesuatu?
    \item Kunci utamanya: kotak awal yang isinya $N=50^{25}$ batu itu \emph{selalu} kehilangan tepat $1$ batu tiap langkah. Jadi kapan dia habis? Dan pas dia habis, lagi ada berapa kotak?
    \item Terakhir, banyaknya kotak yang lagi "hidup" (yang berisi batu) ternyata nyambung sama bilangan segitiga $T_k = \tfrac{k(k+1)}{2}$. Tinggal cari $k$ yang pas.
\end{itemize}


Ayo coba kita jalankan operasinya untuk $N=6$ batu sebagai gambaran (angka menyatakan isi tiap kotak, diurutkan dari yang terbanyak):
\begin{align*}
    (6) \to (5,1) \to (4,2) \to (3,3) \to (3,2,1) \to (3,2,1) \to \cdots
\end{align*}
Perhatikan bahwa kotak awal berisi $6$ batu berkurang $1$ tiap langkah, dan tepat setelah $6$ langkah kotak itu habis. Saat itu banyaknya kotak yang hidup adalah $3$, dan kebetulan $T_3 = \tfrac{3\cdot 4}{2} = 6 = N$.

\newpage
\section{Solusi}
Misalkan $N = 50^{25}$. Pandang \emph{kotak awal} yang mula-mula berisi $N$ batu. Pada setiap operasi, setiap kotak yang tidak kosong kehilangan tepat $1$ batu. Dengan demikian, selama masih ada isinya, kotak awal kehilangan tepat $1$ batu per langkah, sehingga kotak awal menjadi kosong \emph{tepat} setelah operasi ke-$N$. Karena kita memang berhenti setelah operasi ke-$N = 50^{25}$, kita cukup menghitung banyaknya kotak yang hidup (kotak yang berisi batu) tepat pada saat itu.

\medskip
\textbf{Klaim.} Setelah operasi ke-$N$, banyaknya kotak yang hidup adalah bilangan bulat $k$ terbesar yang memenuhi $T_k = \dfrac{k(k+1)}{2} \le N$.

\medskip
Untuk memahami klaim ini, perhatikan bahwa setiap langkah menghapus $1$ batu dari tiap kotak lalu menambahkan satu kotak baru yang isinya sama dengan \emph{banyaknya} kotak yang tadi hidup. Konfigurasi yang stabil (tidak berubah bentuk) adalah tumpukan $ (k, k-1, \dots, 2, 1)$, yang totalnya tepat $T_k$ batu. Untuk $N$ dengan $T_k \le N < T_{k+1}$, proses ini mengendap pada konfigurasi dengan $k$ kotak hidup, dan tepat pada langkah ke-$N$ (saat kotak awal habis) banyaknya kotak hidup sudah bernilai $k$ tersebut.

Jadi $k$ adalah bilangan bulat unik yang memenuhi
\begin{align*}
    \dfrac{k(k+1)}{2} \le N < \dfrac{(k+1)(k+2)}{2},
\end{align*}
yang setara dengan
\begin{align*}
    k(k+1) \le 2N < (k+1)(k+2).
\end{align*}

Sekarang substitusikan $N = 50^{25}$ pada $2N$:
\begin{align*}
    2N = 2\cdot 50^{25} = 100 \cdot 50^{24} = \left(10\cdot 50^{12}\right)^2.
\end{align*}
Misalkan $X = 10\cdot 50^{12}$, sehingga $2N = X^2$ dan pertidaksamaannya menjadi
\begin{align*}
    k(k+1) \le X^2 < (k+1)(k+2).
\end{align*}
Kita uji $k = X-1$:
\begin{align*}
    k(k+1) &= (X-1)X = X^2 - X \le X^2, \\
    (k+1)(k+2) &= X(X+1) = X^2 + X > X^2,
\end{align*}
sehingga kedua pertidaksamaan terpenuhi. Karena $k$ tunggal, haruslah $k = X - 1$.

Dengan demikian, banyaknya kotak yang berisi paling sedikit $1$ batu setelah operasi ke-$50^{25}$ adalah
\begin{align*}
    \boxed{10\cdot 50^{12} - 1}.
\end{align*}
\end{solusi}